

\subsubsection*{Measure deferred tax from the newly assessed book}
The current-tax journal changes the book on which deferred tax will be measured.
In the cumulative example, payable was $27/4$ before the replacement and $35/4$
after it. The declarations used to prepare that current-tax calculation cannot
silently serve as measurements of the later book. The new deferred request
therefore names the exact recorded state and date, declares its jurisdiction and
regime, and supplies fresh bases for every scoped account and current lease
right. It also supplies the item-level bases and recognition decisions that
must reconcile to those totals. An unchanged base may be declared again with
its provenance; the method does not carry that premise forward automatically.

The selected hypothetical model retains the gross, single-recovery-pool,
profit-or-loss calculation described earlier. The complete current population
contains the ordinary asset and liability accounts, inventory subdivisions,
lease rights and remaining loss lots. Fresh measurement and recognition support
is required for every item. Recognized deductions across debt, losses and other
deductible items share the one explicitly supplied recovery capacity. The older
recognition schedule remains historical evidence. Its capacity cannot renew
itself, and a deferred liability cannot silently supply additional capacity.

Recorded allocations must also reconcile to the balances actually carried into
the period. In the worked book, those balances are a deferred tax asset of 8
and a deferred tax liability of 10. With property carrying amount 100 and a
fresh base of 80, the taxable difference is 20. With debt carrying amount 100
and base 60, the deductible difference is 40; suppose fresh support permits
recognition of only 20 against capacity 20. At a declared rate of one quarter,
the desired asset and liability are both 5. The journal reduces the asset by 3,
reduces the liability by 5, and records deferred-tax expense of $-2$.

After the preceding cumulative current-tax example, cash remains $715/2$,
current expense remains 7 and current-tax payable remains $35/4$. Recorded
profit rises from $-19$ to $-17$. Assets become $1361/2$, liabilities $455/4$
and equity including profit $2267/4$; operating cash flow is still $-3$.
All tax-vintage settlements, loss quantities, component states and earlier
findings remain unchanged. The result retains both the fresh recognition
premises and the previous period's schedule.

The capacity constraint matters when loss recognition is added. If the current
loss inventory has 9 remaining and the caller recognizes all 9 as well as the
20 debt deduction, capacity must cover 29. At one quarter, the combined desired
asset is $29/4$. Capacity of only 20 refuses that combination; the same amount
cannot support the debt deduction and then be used again for the loss. A loss
quantity by itself does not establish a deferred asset.

The method compares the entire proposed result when applied, including the
fresh measurements, item allocations, tax vintages and historical recognition
record. Unknown bases or support remain unresolved, and unsupported equity,
other-comprehensive-income and offset routes are refused. Another financial
boundary must preserve these histories while distinguishing old settlements
from the next period's cash movements. Conditional recognition here does not
establish annual completeness or source and accounting admission.
